% Generated mechanically from the frozen desirables.tex target.
% Do not edit this generated file; edit the bound locale source instead.

% OK, start here.
%

\setcounter{chapter}{112}
\chapter{待增内容}


\phantomsection
\label{desirables-section-phantom}



\section{引言}
\label{desirables-section-introduction}

\noindent
这里基本上只是列出我们希望加入 Stacks 项目的内容。由于我们持续向
Stacks 项目添加材料，这份清单总是会稍微落后于项目的当前状态。事实上，
试图列出应当添加的内容也许是个错误，因为这份清单似乎不可能保持最新。

\medskip\noindent
最后更新：2017 年 8 月 31 日，星期四。


\section{约定}
\label{desirables-section-conventions}

\noindent
我们应当增加一章，简要列出文档中使用的约定。
该章节已经存在，见
Conventions，第 \ref{conventions-section-comments} 节，
但那里还可以补充许多内容。尤其有用的是找出“隐藏的”约定和未明说的假设，
并将它们放在该章节中。


\section{位点与拓扑斯}
\label{desirables-section-sites}

\noindent
我们有一章讨论位点与层，见
Sites，第 \ref{sites-section-introduction} 节。
我们有一章讨论环化位点（以及拓扑斯）及其上的模，见
Modules on Sites，第 \ref{sites-modules-section-introduction} 节。
我们还有一章讨论这一背景下的上同调，见
Cohomology on Sites，第 \ref{sites-cohomology-section-introduction} 节。
不过仍可加入更多内容，尤其是上同调章节。


\section{叠}
\label{desirables-section-stacks}

\noindent
我们有一章讨论（抽象的）叠，见
Stacks，第 \ref{stacks-section-introduction} 节。
如果能够做到以下几点就很好：
\begin{enumerate}
\item 改进关于“叠化”的讨论，
\item 给出叠化的例子，
\item 总体上增加更多例子，
\item 改进关于 gerbe 的讨论。
\end{enumerate}
尚未加入的一个结果例子是：给定阿贝尔群层 $\mathcal{F}$
在范畴 $\mathcal{C}$ 上，由 $\mathcal{F}$ 带化的 gerbe 的等价类集合与
$H^2(\mathcal{C}, \mathcal{F})$ 双射。


\section{单纯方法}
\label{desirables-section-simplicial}

\noindent
我们有一章讨论单纯方法，见
Simplicial，第 \ref{simplicial-section-introduction} 节。
这一章需要审阅和改进。关于单纯同伦（也称为组合同伦）与 Kan 复形之间
关系的讨论应当加强。
我们还有一章讨论单纯空间，见
Simplicial Spaces，第 \ref{spaces-simplicial-section-introduction} 节。
该章简要讨论单纯拓扑空间、单纯位点和单纯拓扑斯。
我们可以进一步发展“单纯代数几何”，讨论单纯概形（或单纯代数空间、或
单纯代数叠），并处理几何问题、它们的上同调等。


\section{概形的上同调}
\label{desirables-section-schemes-cohomology}

\noindent
已经有一章讨论拟凝聚层的上同调，见
Cohomology of Schemes，第 \ref{coherent-section-introduction} 节。
我们有一章讨论概形上的拟凝聚层的导出范畴，见
Derived Categories of Schemes，第 \ref{perfect-section-introduction} 节。
我们有一章讨论 Noether 概形的对偶性以及概形态射的相对对偶性，见
Duality for Schemes，第 \ref{duality-section-introduction} 节。
我们还有关于概形的 \'etale 上同调以及概形的结晶上同调的章节。不过这些
章节中的大部分材料都非常基础，其中还有许多内容可以、也应当补充。


\section{Schlessinger 风格的形变理论（\`a la Schlessinger）}
\label{desirables-section-deformation-schlessinger}

\noindent
我们有一章讨论这方面的材料，见
Formal Deformation Theory，第 \ref{formal-defos-section-introduction} 节。
我们有一章讨论一般理论的例子，见
Deformation Problems，第 \ref{examples-defos-section-introduction} 节。
我们还有一章，见
Deformation Theory，第 \ref{defos-section-introduction} 节，
其中讨论环（以及模）的形变、环化空间（以及模层）的形变、以及环化拓扑斯
（以及模层）的形变。
本章使用朴素余切复形描述障碍、一阶形变和无穷小自同构。这些材料后来已
在模叠的代数性问题中得到一些应用。
另有一章讨论完整余切复形，见
Cotangent，第 \ref{cotangent-section-introduction} 节。


\section{代数叠的定义}
\label{desirables-section-definition-algebraic-stacks}

\noindent
代数叠是在带有 fppf 拓扑的概形范畴上的群胚中的叠，其对角态射可由代数
空间表示，并且是某个概形到它的满射光滑态射的目标。
见 Algebraic Stacks，第 \ref{algebraic-section-algebraic-stacks} 节。
“Deligne--Mumford 叠”是满足下述条件的代数叠：存在一个概形以及从该概形
到它的满射 \'etale 态射，正如 Deligne 和 Mumford 的论文 \cite{DM} 中所述；
见 Algebraic Stacks，定义 \ref{algebraic-definition-deligne-mumford}。
我们将“Artin 叠”一词保留给 Artin 的论文中所讨论的叠，见
\cite{ArtinI}、\cite{ArtinII} 和 \cite{ArtinVersal}。
一种可能的定义是：Artin 叠是局部 Noether 概形 $S$ 上的代数叠
$\mathcal{X}$，使得 $\mathcal{X} \to S$ 局部有限型
\footnote{也就是说，这些恰好是满足 Artin 公理 [-1]、[0]、[1]、[2]、[3]、[4]、[5]
的 $S$ 上代数叠，见 Artin's Axioms，第 \ref{artin-section-axioms} 节。}。


\section{概形、代数空间与代数叠的例子}
\label{desirables-section-examples-stacks}

\noindent
Stacks 项目目前有两章讨论模叠及其性质，见
Moduli Stacks，第 \ref{moduli-section-introduction} 节，以及
Moduli of Curves，第 \ref{moduli-curves-section-introduction} 节。
随着时间推移，我们打算加入更多例子，例如：
\begin{enumerate}
\item $\mathcal{A}_g$，即亏格为 $g$ 的主极化阿贝尔概形，
\item $\mathcal{A}_1 = \mathcal{M}_{1, 1}$，即带 $1$ 个标记点的光滑射影亏格 $1$ 曲线，
\item $\mathcal{M}_{g, n}$，即带有 $n$ 个两两不同且带标号点的光滑射影亏格 $g$ 曲线，
\item $\overline{\mathcal{M}}_{g, n}$，即带 $n$ 个标记点的稳定结点射影亏格 $g$ 曲线，
\item $\SheafHom_S(\mathcal{X}, \mathcal{Y})$，态射的模空间
（对叠 $\mathcal{X}$、$\mathcal{Y}$ 以及基概形 $S$ 施加适当条件），
\item $\textit{Bun}_G(X) = \SheafHom_S(X, BG)$，几何 Langlands 纲领中的 $G$-丛叠
（对概形 $X$、群概形 $G$ 以及基概形 $S$ 施加适当条件），
\item $\Picardstack_{\mathcal{X}/S}$，即与基概形（或基空间）上的代数叠关联的 Picard 叠。
\end{enumerate}
更一般地说，Stacks 项目在具有几何意义的例子方面仍然有所欠缺。


\section{代数叠的性质}
\label{desirables-section-stacks-properties}

\noindent
这也许是比较容易开展的项目之一，因为大部分基础理论现在已经具备。
当然，这些实际上是叠的态射的性质。我们可以定义奇点（模去光滑因子）等；
还可以证明连通正规叠是不可约的，等等。


\section{代数叠的光滑 \'etale 位点}
\label{desirables-section-lisse-etale}

\noindent
这一内容已在 Cohomology of Stacks，第 \ref{stacks-cohomology-section-lisse-etale} 节中介绍。
为了说明它对于代数叠的 $1$-态射并非函子性，Examples，第
\ref{examples-section-lisse-etale-not-functorial} 节中讨论了一个例子。
当然还可以说得更多，但事实证明，尽可能使用“大” \'etale 位点来证明结果
非常有用。


\section{你一直想知道却不敢问的事情}
\label{desirables-section-stacks-fun-lemmas}

\noindent
会有许多引理被反复使用：它们很有用，却没有在文献中得到明确提及，或者
很难找到参考文献。这里就是一个技巧袋。

\medskip\noindent
例：给定概形中的两个群胚 $R\Rightarrow U$ 和
$R' \Rightarrow U'$，仅用概形中的群胚来描述一个
$1$-态射 $[U/R] \to [U'/R']$ 到底意味着什么。


\section{叠上的拟凝聚层}
\label{desirables-section-quasi-coherent}

\noindent
这些内容在 Cohomology of Stacks 章节中定义和讨论，见
第 \ref{stacks-cohomology-section-introduction} 节。
模的导出范畴在 Derived Categories of Stacks 章节中讨论，见
第 \ref{stacks-perfect-section-introduction} 节。
这些章节还可以补充许多内容。


\section{平坦与光滑}
\label{desirables-section-flat-smooth}

\noindent
Artin 定理说明，从概形出发的平坦满射可以替代满射光滑条件。
这一结果现在已作为 Criteria for Representability，第
\ref{criteria-theorem-bootstrap} 条定理提供。


\section{Artin 表示定理}
\label{desirables-section-representability}

\noindent
这一内容在 Artin's Axioms 章节中讨论，见
第 \ref{artin-section-introduction} 节。
我们还有一个应用，见
Quot，第 \ref{quot-theorem-coherent-algebraic} 条定理。
还应当有更多应用，并且该章节本身也需要清理。


\section{DM 叠由概形有限覆盖}
\label{desirables-section-dm-finite-cover}

\noindent
对于代数空间，我们已经有相应结果，见 Limits of Spaces，第
\ref{spaces-limits-section-finite-cover} 节。
缺少的是 DM 叠和拟 DM 叠的结果。


\section{Martin Olsson 关于真性的论文}
\label{desirables-section-proper-parametrization}

\noindent
该论文证明两种真性的概念相同。第一部分现在以代数叠的 Chow 引理形式
提供，见 More on Morphisms of Stacks，第
\ref{stacks-more-morphisms-theorem-chow-finite-type} 条定理。
作为推论，在某些情形下检验代数叠的赋值判据时，只需使用 DVR，见
More on Morphisms of Stacks，第
\ref{stacks-more-morphisms-section-Noetherian-valuative-criterion} 节。


\section{相干层的真推前}
\label{desirables-section-proper-pushforward}

\noindent
现在我们已经有了代数叠的 Chow 引理，可以开始研究这一内容，见上一节。


\section{Keel 与 Mori}
\label{desirables-section-keel-mori}

\noindent
见 \cite{K-M}。他们的结果已经加入 More on Morphisms of Stacks，第
\ref{stacks-more-morphisms-section-Keel-Mori} 节。


\section{在此处添加更多内容}
\label{desirables-section-add-more}

\noindent
其实，我们根本不应当把这份清单作为 Stacks 项目本身的一部分开始！
其他地方有一份待办清单，更容易更新。
